% Integrated release technical supplement, revised 15 October 2026.
% Quantitative macros are generated from the rebuilt rc68 products.
% Official Copernicus class remains unmodified.
\documentclass[essd]{copernicus}

% siunitx AND THE SHARED MACROS, on the same terms as paper.tex - absent until 2026-10-07,
% which meant this file did not compile at all. It uses \SI and \num throughout and no
% target built it (see the paper-v3-supplement rule in the Makefile), so the two defects
% hid each other: nothing compiled the file, so nobody saw that it could not compile.
%
% The ordering is not cosmetic and is copied from paper.tex deliberately: the class loads
% its own fontenc/inputenc, graphicx, amsmath, natbib, hyperref, tabularx, multirow, float
% and caption, and loading any of those again breaks a Copernicus build; siunitx must
% precede the macros because they call \DeclareSIUnit; and the class defines a robust
% \unit that siunitx v3 refuses to define over, so the class's copy is released first.
\let\unit\relax
\usepackage{siunitx}
% ===========================================================================
% paper_common/macros.tex - the shorthand every manuscript shares.
%
% Split from preamble.tex in the ESSD re-targeting: the national article now compiles
% under the Copernicus class, which loads its own packages and forbids most of what
% preamble.tex sets (geometry-like options, a second hyperref, placeins), but it needs
% exactly the same shorthand as the sector reports. The macros therefore live here, on
% their own, and preamble.tex inputs this file so that every other manuscript keeps
% compiling unchanged. A macro that renders differently in two documents is the defect
% this file exists to prevent; define it once.
%
% Requires: siunitx (for \DeclareSIUnit and \si) - the caller loads it.
% ===========================================================================

% Gas and pollutant names.
\newcommand{\COtwo}{CO\textsubscript{2}}
\newcommand{\COe}{CO\textsubscript{2}e}
\newcommand{\CHfour}{CH\textsubscript{4}}
\newcommand{\NtwoO}{N\textsubscript{2}O}
\newcommand{\NHthree}{NH\textsubscript{3}}
\newcommand{\NOx}{NO\textsubscript{x}}
\newcommand{\SOtwo}{SO\textsubscript{2}}
\newcommand{\PM}{PM\textsubscript{2.5}}
\newcommand{\MCF}{MCF}
\newcommand{\Mt}{\text{Mt~\COe}}
\newcommand{\tCha}{\si{\tonne\per\hectare}}
\newcommand{\AFOLU}{AFOLU}
\newcommand{\GLW}{GLW\,4}
\newcommand{\EFone}{EF\textsubscript{1}}

% siunitx: allow Mt, Tg, ha, head, yr, dB
\DeclareSIUnit{\head}{head}
\DeclareSIUnit{\yr}{yr}
\DeclareSIUnit{\dB}{dB}

% A file name is the one long unbreakable token these documents are full of -
% \texttt{waste\_air\_pollutants\_2024.csv} is wider than a two-column measure and TeX prints
% it past the margin rather than breaking it. Allowing a break AFTER each escaped underscore
% costs nothing and inserts no hyphen, which matters: a hyphenated file name is a wrong file
% name. This is one line because every such name in every document is written with \_.
\let\pkinvunderscore\_
\renewcommand{\_}{\pkinvunderscore\allowbreak}


% Generated from released native gas and pollutant products; do not hand edit.
\newcommand{\PakEmisRelease}{1.0.0rc68}
\newcommand{\CoreEndYear}{2025}
\newcommand{\PakRefCOtwo}{161}
\newcommand{\PakRefCHfour}{8621}
\newcommand{\PakRefNtwoO}{153.4}
\newcommand{\PakRefARfive}{444.6}
\newcommand{\PakRefARsix}{437.8}
\newcommand{\PakRefPMfine}{422.8}
\newcommand{\PakRefBC}{40.44}
\newcommand{\PakRefNOx}{816.6}
\newcommand{\PakRefSOtwo}{213.6}
\newcommand{\PakRefCO}{5991}
\newcommand{\PakRefNMVOC}{1187}
\newcommand{\PakRefNHthree}{2670}
\newcommand{\PakLatestCOtwo}{167.6}
\newcommand{\PakLatestCHfour}{8848}
\newcommand{\PakLatestNtwoO}{152.6}
\newcommand{\PakLatestARfive}{457.2}
\newcommand{\PakLatestARsix}{450.2}
\newcommand{\PakLatestPMfine}{422.7}
\newcommand{\PakLatestBC}{46.64}
\newcommand{\PakLatestNOx}{833.1}
\newcommand{\PakLatestSOtwo}{368.9}
\newcommand{\PakLatestCO}{5999}
\newcommand{\PakLatestNMVOC}{1226}
\newcommand{\PakLatestNHthree}{2690}
\newcommand{\PakFirstCOtwo}{173}
\newcommand{\PakFirstCHfour}{6895}
\newcommand{\PakFirstNtwoO}{113.9}
\newcommand{\PakFirstARfive}{397.8}
\newcommand{\PakFirstARsix}{392.5}
\newcommand{\PakFirstPMfine}{340.8}
\newcommand{\PakFirstBC}{40.94}
\newcommand{\PakFirstNOx}{705.1}
\newcommand{\PakFirstSOtwo}{671.2}
\newcommand{\PakFirstCO}{4700}
\newcommand{\PakFirstNMVOC}{942.1}
\newcommand{\PakFirstNHthree}{1847}
\newcommand{\PakChangeARfive}{14.9}
\newcommand{\PakChangePMfine}{24.0}
\newcommand{\PakChangeBC}{13.9}
\newcommand{\PakChangeNOx}{18.2}
\newcommand{\PakChangeSOtwo}{-45.0}
\newcommand{\PakChangeCO}{27.6}
\newcommand{\PakChangeNMVOC}{30.2}
\newcommand{\PakChangeNHthree}{45.7}
\newcommand{\PakRefBCLow}{27.42}
\newcommand{\PakRefBCMedian}{39.04}
\newcommand{\PakRefBCHigh}{62.03}
\newcommand{\PakRefCOLow}{4777}
\newcommand{\PakRefCOMedian}{5942}
\newcommand{\PakRefCOHigh}{7497}
\newcommand{\PakRefNHthreeLow}{2297}
\newcommand{\PakRefNHthreeMedian}{2657}
\newcommand{\PakRefNHthreeHigh}{3069}
\newcommand{\PakRefNMVOCLow}{944.5}
\newcommand{\PakRefNMVOCMedian}{1176}
\newcommand{\PakRefNMVOCHigh}{1507}
\newcommand{\PakRefNOxLow}{660.6}
\newcommand{\PakRefNOxMedian}{806.4}
\newcommand{\PakRefNOxHigh}{1029}
\newcommand{\PakRefPMfineLow}{311.2}
\newcommand{\PakRefPMfineMedian}{416.8}
\newcommand{\PakRefPMfineHigh}{567}
\newcommand{\PakRefSOtwoLow}{163}
\newcommand{\PakRefSOtwoMedian}{208.1}
\newcommand{\PakRefSOtwoHigh}{293}
\newcommand{\PakLatestBCLow}{31.3}
\newcommand{\PakLatestBCMedian}{45.01}
\newcommand{\PakLatestBCHigh}{71.23}
\newcommand{\PakLatestCOLow}{4796}
\newcommand{\PakLatestCOMedian}{5947}
\newcommand{\PakLatestCOHigh}{7463}
\newcommand{\PakLatestNHthreeLow}{2333}
\newcommand{\PakLatestNHthreeMedian}{2680}
\newcommand{\PakLatestNHthreeHigh}{3061}
\newcommand{\PakLatestNMVOCLow}{976.4}
\newcommand{\PakLatestNMVOCMedian}{1215}
\newcommand{\PakLatestNMVOCHigh}{1556}
\newcommand{\PakLatestNOxLow}{679.5}
\newcommand{\PakLatestNOxMedian}{823.3}
\newcommand{\PakLatestNOxHigh}{1040}
\newcommand{\PakLatestPMfineLow}{311.7}
\newcommand{\PakLatestPMfineMedian}{416.6}
\newcommand{\PakLatestPMfineHigh}{564.9}
\newcommand{\PakLatestSOtwoLow}{221.8}
\newcommand{\PakLatestSOtwoMedian}{348.7}
\newcommand{\PakLatestSOtwoHigh}{616.7}
\newcommand{\PakRefARfiveLow}{399.9}
\newcommand{\PakRefARfiveMedian}{445.9}
\newcommand{\PakRefARfiveHigh}{503.2}
\newcommand{\PakLatestARfiveLow}{411.4}
\newcommand{\PakLatestARfiveMedian}{458.6}
\newcommand{\PakLatestARfiveHigh}{517.3}
\newcommand{\PakDungCaseBC}{61.89}
\newcommand{\PakDungCasePMfine}{772}
\newcommand{\PakRefNPartitionNHthreeChangeGg}{24.65}
\newcommand{\PakRefNPartitionNOxChangeGg}{-149.8}
\newcommand{\PakLatestNPartitionNHthreeChangeGg}{26.06}
\newcommand{\PakLatestNPartitionNOxChangeGg}{-158.1}
\newcommand{\PakRefNPartitionGHGChange}{0.04426}
\newcommand{\PakRefNPartitionARsixChange}{0.04559}
\newcommand{\PakLatestNPartitionGHGChange}{0.04736}
\newcommand{\PakLatestNPartitionARsixChange}{0.04879}
\newcommand{\GHGReferenceTotal}{\PakRefARfive}
\newcommand{\GHGLatestTotal}{\PakLatestARfive}
\newcommand{\GHGReferenceLow}{\PakRefARfiveLow}
\newcommand{\GHGReferenceHigh}{\PakRefARfiveHigh}
\newcommand{\GHGLatestLow}{\PakLatestARfiveLow}
\newcommand{\GHGLatestHigh}{\PakLatestARfiveHigh}
\newcommand{\PakNPartitionNOxChangeGg}{\PakRefNPartitionNOxChangeGg}
\newcommand{\PakNPartitionNHthreeChangeGg}{\PakRefNPartitionNHthreeChangeGg}
\newcommand{\PakNPartitionGHGChange}{\PakRefNPartitionGHGChange}
\newcommand{\PakNPartitionARsixChange}{\PakRefNPartitionARsixChange}
\newcommand{\AirReferencePM}{\PakRefPMfine}
\newcommand{\AirReferenceBC}{\PakRefBC}
\newcommand{\AirReferenceNOx}{\PakRefNOx}
\newcommand{\AirReferenceSOtwo}{\PakRefSOtwo}
\newcommand{\AirReferenceCO}{\PakRefCO}
\newcommand{\AirReferenceNMVOC}{\PakRefNMVOC}
\newcommand{\AirReferenceNHthree}{\PakRefNHthree}
\newcommand{\AirLatestPM}{\PakLatestPMfine}
\newcommand{\AirLatestBC}{\PakLatestBC}
\newcommand{\AirLatestNOx}{\PakLatestNOx}
\newcommand{\AirLatestSOtwo}{\PakLatestSOtwo}
\newcommand{\AirLatestCO}{\PakLatestCO}
\newcommand{\AirLatestNMVOC}{\PakLatestNMVOC}
\newcommand{\AirLatestNHthree}{\PakLatestNHthree}


\begin{document}
\title{Technical documentation for PakEmis \PakEmisRelease: greenhouse gas and air pollutant emissions, July 2015--June 2025}
\Author[1][rehan@pakairquality.com]{Rehan}{Ahmad}
\affil[1]{Pakistan Air Quality Initiative (PAQI), Karachi, Pakistan}
\runningtitle{PakEmis integrated technical documentation}
\runningauthor{R. Ahmad}
\received{}\pubdiscuss{}\revised{}\accepted{}\published{}
\firstpage{1}\maketitle

\begin{abstract}
This documentation describes the shared activity accounting, chemical conventions, source-specific calculations and reproduction interfaces underlying the national greenhouse-gas and air-pollutant dataset. The common core comprises ten fiscal years, July 2015--June 2025 (FY2015--16 to FY2024--25); FY2023--24 is the detailed reference year. Source-specific provisional estimates for FY2025--26 are kept separate. This document accompanies the national data-description article and the maintained sector notes and does not define another inventory total. Numerical summaries are read from the same released products as the main article. Appendix~\ref{app:products} lists every canonical data product.
\end{abstract}

\section{Scope, period and record conventions}
The native greenhouse gases are CO$_2$, CH$_4$ and N$_2$O. The core air panel is primary fine particulate matter (PM$_{2.5}$), black carbon (BC), nitrogen oxides (NO$_x$), sulfur dioxide (SO$_2$), carbon monoxide (CO), non-methane volatile organic compounds (NMVOC) and ammonia (NH$_3$). PM$_{10}$, total suspended particles (TSP), organic carbon (OC) and other species have narrower source coverage. Land use, land-use change and forestry are outside the inventory. Fluorinated gases, coal-mining fugitive emissions, several industrial processes and some treatment and combustion sources are listed as unestimated. A product that spans all modules is therefore not a complete category--species matrix.

A fiscal year runs from July to June and is identified by its start and end years; the FY2024--25 core endpoint is June 2025, not calendar 2025. A monthly publication through June 2026 can support a full FY2025--26 physical series when all required months and variables are present, whereas the Economic Survey 2025--26 reports several series for July--March only. The availability table records edition, covered months, revision status, variable, unit, held/public location and usability before a source is admitted to an extension.

The annual records identify the source, process, reporting category and calculation module separately. They retain native gas/species, mass convention, activity quantity and unit, activity period and provenance, factor unit and applicability, technology, control state, configuration, estimate status, temporal/spatial support and uncertainty model. Estimated zero, not applicable and not estimated are distinct statuses. Source coverage is supplied by process, species and year; configuration alternatives are never implicitly added to the reference.

NO$_x$ is NO$_2$-equivalent mass, with NO-N converted by $46/14$. NH$_3$-N is converted to NH$_3$ by $17/14$, and N$_2$O-N to N$_2$O by $44/28$. Native pollutant masses never use a global warming potential. The separate AR5 and AR6 aggregates identify the selected fossil/biogenic methane convention \citep{IPCCAR5,IPCCAR6}. BC and OC cannot be added to compatible PM$_{2.5}$ totals; nor can PM$_{2.5}$, PM$_{10}$ and TSP be summed across nested definitions. Precursors are not combined into emitted ozone or total ambient PM$_{2.5}$.

\section{Source ingestion and shared activity}
\subsection{Monthly minerals and source controls}
Monthly mineral production from the Pakistan Bureau of Statistics (PBS) is read at its printed mineral, year and provincial grouping. The printed year block, not the workbook name, identifies the year; worksheets are selected by table schema and provenance; and a province row keeps the mineral group whose label and printed total enclose it. A printed dash or an illegible cell is kept as a missing value rather than a zero, and a province absent from the source is not added.

Exported records are compared with independently transcribed source quantities: province--month entries, published national months, compatible annual totals and the expected producing provinces. A twelve-month profile is classed as national observed timing only when both its months and its producing-area coverage support that reading; otherwise the profile and its evidence class name the substitute. Mineral production supports the timing of the corresponding upstream oil and gas segments, while coal mine output does not establish when kilns or power stations burn coal. Because normalisation removes a uniform duplication, an unchanged profile cannot validate an exported activity table, so source controls, exported quantities and profiles are checked separately.

\subsection{Coal mass, energy and consumer allocation}
Station coal consumption is reported in kt (1~kt = 1~Gg) by the National Electric Power Regulatory Authority (NEPRA) \citep{NEPRA2024}. For station $i$ and rank $r$,
\begin{equation}
 Q_{i,t}=M_{i,t}\,NCV_{r(i),t},\qquad
 E_{\mathrm{CO_2},t}=\sum_i Q_{i,t}EF_{\mathrm{CO_2},r(i)},
 \label{eq:station_coal}
\end{equation}
where $Q$ is TJ, $M$ is kt and NCV is TJ\,Gg$^{-1}$. Thus the blend in kg\,CO$_2$\,TJ$^{-1}$ is weighted by heat, not by generation. The reference retains the documented rank-default NCVs. Reported coal used is distinct from purchases, receipts, stock movements and generation. Testing-period consumption is included where it is actually reported.

When a commissioning lane has observed generation but no printed mass, the reconstruction retains an interval derived from the documented station-intensity evidence. Its midpoint defines the conditional reference calculation; the endpoints carry no assigned probability or confidence level. A missing latest-year station transcription causes a failure rather than assigning its entire generation to a reconstructed lane. Older aggregate coal energy remains a comparator, and a year without a published aggregate energy value remains absent in that comparator.

The national coal account retains production, imports, exports and stock change where observed, and the consumer columns for public power, households, brick kilns and cement/other industry. Captive and K-Electric-system coal have no printed fuel quantity and remain, unsplit, inside the published cement and other industry column. If exports or stocks are not observed, the resulting discrepancy is not labelled a measured physical residual. The source records' coal origin, rank, moisture or calorific basis remain distinct.

The adopted industrial publication vintage is selected explicitly. For FY2024--25 the full-year table in the later Economic Survey assigns brick kilns and cement and other industry differently from the earlier provisional supplement, with the same combined activity. Both versions remain traceable. Substituting station consumption for a publisher's power allocation does not establish which non-power consumer used the difference. The unresolved quantity remains in the account and is not automatically emitted as kiln fuel.

Dated kiln counts are evaluated with the stated fuel rate and operating-season assumptions. The Punjab fleet cases retain their different observation dates and are scaled to the national total by the Punjab share of counted kilns; each count is compared with the booked allocation only in the fiscal year in which it was made, and the cases are not averaged. Negative supply-minus-use quantities under a case are retained as contradictions or unresolved source differences, not rounded away. Source-specific GHG and pollutant effects of these cases are reported separately from the conditional activity distribution.

\subsection{Delivered gas, processes and transport fuel}
The gas ledger distinguishes production, processing, utility delivery, chemical feedstock, process fuel and other combustion. Fertiliser feedstock is separated from combustion fuel where the source supports it. It is not included in the post-meter combustion-leakage base merely because a utility calls both quantities fertiliser deliveries. Methane releases are estimated with segment-specific activities and factors; national unaccounted-for gas is not used as a measured methane loss.

The shared gas deduction acts on the combustion activity to which the post-meter factor applies. A draw assigning additional methane release correspondingly changes the combusted fraction. This coupling prevents the same gas from being fully burned and lost as methane. It does not resolve all carbon differences between national gas statistics and source-specific process production.

Petroleum product sales retain customer-class and fuel distinctions. Retail diesel is partitioned between road use and off-road machinery; petrol attributed to residential generators is deducted from the quantity otherwise available to road transport. Operator fuel distinguishes railway traction from non-traction activity. Domestic aviation is separated from international bunkers and landing/take-off from cruise where the evidence permits. Electrical energy bought from another generator is not counted as a second combustion source. K-Electric's fuel-specific station record and purchased generation are therefore separate boundary questions. Its FY2024--25 gas and regasified liquefied natural gas (RLNG) coverage case converts gross-heat volumes to the station net-calorific basis before comparison or emissions calculation; oil and coal use their reported physical fuel quantities.

\section{Agricultural nitrogen and household configurations}
\subsection{Census-consistent historical activity}
Historical livestock levels are anchored to the adopted 2024 census \citep{PBSAgri2026}. For species $a$ and proxy stock $P_{a,t}$,
\begin{equation}
 H_{a,t}=H_{a,2024}\frac{P_{a,t}}{P_{a,2024}}.
 \label{eq:livestock_bridge}
\end{equation}
The proxy supplies relative changes, not a second census level. It is the FAOSTAT livestock series, which reproduces the Economic Survey's projection from 1996--2006 inter-census growth rates, so the year-to-year change of every species, including the FY2024--25 growth, is a projected rate rather than an annual count; the 2024 census differs from that projection most for sheep. Beyond the last held proxy year, a species-specific geometric rate over the declared recent proxy window is applied to the census anchor, which is modelled extrapolation. The phase sensitivity uses the same bridge when assessing the six-month stock-to-fiscal alignment.

Crops use the full-year FY2024--25 tables of the Pakistan Economic Survey 2025--26 and historical fiscal editions; fertiliser product composition and offtake use the full-year Annual Fertilizer Review 2024--25 of the National Fertilizer Development Centre (NFDC) and historical editions \citep{MoF2024,MoF2026,NFDC2024,NFDC2025}. Cotton lint bales are converted to the seed-cotton mass required by the residue equations with the documented conversion. Crop production, area and fertiliser product composition retain distinct uncertainty and applicability. Fertiliser offtake is a supply quantity; it is not an observed application date or a direct measurement of field nitrogen loss.

\subsection{Conserved management and field pathways}
Nitrogen excretion for species $a$ is
\begin{equation}
 N_{\mathrm{ex},a}=H_a r_{N,a} M_a\,365/1000,
 \label{eq:excretion}
\end{equation}
where $M_a$ is typical animal mass and $r_{N,a}$ is the daily nitrogen excretion rate per unit animal mass on the source's stated basis. The total is partitioned into grazing deposition, managed manure and diverted material. Destination assumptions do not change this partition into evidence that all diverted dung was combusted.

In the reference TAN formulation, incoming TAN is the species-specific fraction of the relevant nitrogen pool. At housing, yard or storage stage $j$,
\begin{equation}
 \begin{aligned}
 L_{\mathrm{NH_3-N},j}&=TAN_j f_j,\\
 TAN_{j+1}&=TAN_j-L_{\mathrm{NH_3-N},j}-L_{\mathrm{NO_x-N},j}.
 \end{aligned}
 \label{eq:tan}
\end{equation}
The adopted NH$_3$-N share of mixed volatilisation is 0.95, so the associated NO$_x$-N loss in management and grazing is $L_{\mathrm{NH_3-N}}(1-0.95)/0.95$ in the TAN formulation. This chemical attribution is not a measured Pakistani management or grazing factor. Other management losses remove TAN and organic nitrogen in proportion to their surviving pools where no evidence establishes another partition. Nitrogen reaching application is
\begin{equation}
 \begin{aligned}
 N_{\mathrm{applied}}={}&N_{\mathrm{managed}}-L_{\mathrm{NH_3-N}}-L_{\mathrm{NO_x-N}}\\
 &-L_{\mathrm{leach}}-L_{\mathrm{N_2}}-L_{\mathrm{N_2O-N}}.
 \end{aligned}
 \label{eq:fon}
\end{equation}
Each fate is accounted once and closure is tested for every species/system and every uncertainty realisation.

Direct soil N$_2$O uses gross applied nitrogen on its applicable IPCC basis. Application volatilisation drives subsequent indirect N$_2$O and is not deducted again before direct N$_2$O. Grazing retains its distinct deposited nitrogen and factors. Synthetic NH$_3$ follows product-weighted nitrogen offtake and NH$_3$-N factors; the associated NO$_x$-N term is retained when constructing the indirect N$_2$O activity. Product factors already expressed as NH$_3$-N are not multiplied by an additional NH$_3$ share.

The alternative IPCC formulation applies its complete pre-field and field volatilisation fractions to the corresponding pools. The reference upland, flooded-rice and grazing factor choices retain their source-defined applicability \citep{IPCC2019}. Synthetic nitrogen is partitioned to rice where supported; organic and residue inputs do not have a fully observed crop and irrigation allocation. Factor and allocation sensitivities are published, and agreement with FAOSTAT \citep{FAOSTAT2025} is not used to select a factor.

A separate chemical-attribution diagnostic replaces the mixed-loss NH$_3$-N share by 1.0 and reruns manure fate, field inputs, NH$_3$, NO$_x$ and native greenhouse gases for each of the four configurations and all core years. In the TAN case this removes inferred management and grazing NO$_x$-N, changes the surviving manure nitrogen and thereby changes later NH$_3$ and N$_2$O. Separately factored applied-nitrogen and burning NO$_x$ are retained. In the IPCC formulation the same diagnostic repartitions fixed mixed volatilisation losses, including application, while retaining their total nitrogen loss. The diagnostic is neither a measured alternative nor a plausible bound and has no probability. Air-pollutant masses and native-gas/AR5/AR6 consequences occupy separately typed tables; other sources are held at their configuration's central values.

\subsection{Survey-based and diverted-dung household activity}
The household reference estimates fuel use from quantities in the Household Integrated Economic Survey (HIES), household counts from the 2023 census and supported fuel sales \citep{HIES2019,PBSHIES2026,PBSCensus2023}. Interpolation and extrapolation between surveys are labelled. The alternative burns all the material diverted from manure management in households; it is a destination assumption. Under that alternative the uncertain herd drives both diverted dung and household emissions in the same realisation. Volatile solids, dry solids including ash, air-dry fuel and moisture remain distinct mass bases.

\begin{table}[t]
\caption{Structural configurations. The machine identifiers are join keys; the reference is the survey/TAN combination. No probabilities are assigned.}
\label{tab:configurations}
\begin{tabularx}{\linewidth}{@{}lXX@{}}
\tophline
Identifier & Household dung & Nitrogen formulation \\
\middlehline
\texttt{S\_TAN} & Survey fuel use & TAN chain \\
\texttt{L\_TAN} & All diverted dung combusted & TAN chain \\
\texttt{S\_IPCC} & Survey fuel use & Complete IPCC fractions \\
\texttt{L\_IPCC} & All diverted dung combusted & Complete IPCC fractions \\
\bottomhline
\end{tabularx}
\end{table}

Residential pollutant factors represent gross appliance emissions. No nationally justified outdoor escape fraction is assumed. Household emissions should therefore not be interpreted as outdoor atmospheric-model releases without an additional boundary treatment. Unresolved dung destinations remain in a residual material account, with no invented combustion or return-to-soil pathway.

\section{Air-pollutant source equations and coverage}
\subsection{Stationary combustion and processes}
Activity is shared with the greenhouse-gas account, but pollutant factors retain their applicable boiler, furnace, engine, turbine, kiln or process boundary. GHG cement process and fuel emissions remain separate; the cement air factor represents the integrated kiln stack and is applied once to clinker output. Other industrial process and combustion boundaries retain separate source classes. Gas combustion uses the quantity remaining after the shared post-meter leakage deduction. A controlled factor receives no second removal efficiency. A plant licence or EIA can establish installed or designed equipment without establishing actual operation. The reference admitted power oil/diesel SO$_2$ remains distinct from coal screening and other unresolved plant pollutants.

Sulfur-based calculations keep fuel sulfur, retention and removal as separate quantities. Brick-kiln SO$_2$ is calculated from fuel sulfur only, with an assumed retention of 0.25; sulfur from the clay is not estimated. The cement-kiln factor combines fuel and raw-meal sulfur with in-kiln retention and cannot be separated into those parts. Limited kiln-fuel samples constrain particular fuels and dates; they are not a nationally representative sampling distribution. Uncontrolled combustion endpoints are conditional calculations; they are not treated as bounds on actual emissions.

Cement, ammonia, steel, soda ash and other process rows retain only their admitted source boundaries. Unestimated urea finishing, nitric acid, lime, glass or other industrial sources remain in the coverage register. Flaring is separate from venting; the same feed or fuel flow is not added in energy, IPPU and fugitives as three independent inputs.

\subsection{Road, rail and aviation}
Road air emissions comprise the supported exhaust, evaporative, tyre/brake and road-surface wear calculations. Exhaust factors include the guidebook's average hot/cold treatment; the dataset does not supply an independently calibrated Pakistan cold-start model. Registration survival, vehicle age and class/fuel proxies constrain the transferred factors. CNG chemistry and vehicle-class substitutions are labelled where retained. Road resuspension is excluded, and road estimates do not silently include off-road machinery.

Fuel-sulfur policy dates define specification scenarios rather than pump-level compliance observations. Technology and factor endpoints are propagated into pollutant consequences independently of the fuel/CO$_2$ closure. Non-exhaust wear follows distance/vehicle activity rather than treating its emissions as a fixed combustion factor. Rail factors follow operator traction fuel and their technology assumptions; domestic aviation retains its source-specific activity and phase boundaries.

Monthly resident-airline jet-fuel purchases are a proxy for domestic-flight activity: the purchaser and flight-geography boundaries differ. Aviation therefore carries proxy timing support rather than observed monthly domestic movements. The common temporal ensemble's customer-sales hindcast describes dispersion within that proxy profile; it does not quantify the unresolved transfer from customer uplift to domestic-flight timing. National petrol distribution similarly treats road-customer sales as a timing proxy for a broader activity account.

\subsection{Agricultural and waste burning}
Crop burning uses production, residue ratios, burned fractions, combustion completeness and species factors on their documented mass bases. Residue returned to soil, diverted fuel and field burning are accounted on compatible definitions. A source-specific factor does not imply that the burned fraction was annually observed.

Waste generation per person follows the eight-city national survey \citep{PakEPA1996} and is partitioned into collected disposal and uncollected material. The latter is partitioned into burning and unresolved fate. The first-order-decay state carries the disposal history forward, rather than treating an annual population update as an observation of annual waste-management change \citep{IPCC2006,Kaza2018,JMP2023}. Open-burning CH$_4$ and N$_2$O retain their source-specific wet/dry bases. The pollutant calculation uses the same burned dry mass as the native waste uncertainty replay. Waste origin geography does not establish the location of disposal, treatment or burning.

\section{Monthly and geographic products}
For an allocation with annual mass $E_{s,p,t}$ and monthly weights $w_{s,p,t,m}$, following the structure of sector-specific temporal profiles \citep{Crippa2020},
\begin{equation}
 E_{s,p,t,m}=E_{s,p,t}w_{s,p,t,m},\qquad \sum_m w_{s,p,t,m}=1.
 \label{eq:monthly}
\end{equation}
The weights carry their own source and evidence classification. A twelve-month profile with incomplete producing-area coverage is not national observed activity. A flat constant-daily-rate profile weights actual month lengths and leap days, expressing unresolved seasonality. Fertiliser sales, field application, manure stages, evaporation and combustion can require different profiles despite their shared annual input.
FY2024--25 crop-residue burning has supported annual activity but no held own-year VIIRS fire windows. Its monthly mass uses the labelled elapsed-days fallback and its cropland placement uses a substitute surface; neither is an observed fire pattern for that year.
The held crop district panel stops at FY2023--24. Crop-driven FY2024--25 district weights therefore reuse its reference shape and are separately marked as a forward proxy. Residential and waste district weights driven by population projections are marked as modelled demographic patterns, even when the household survey strata are from FY2024--25.

Annual administrative and grid products preserve the distinction between reporting and physical geography. Facility commissioning/closure periods are retained, and historical allocations state whether they use own-year evidence, a fixed proxy, a reference pattern or a substitute. This prevents a repeated location surface from being interpreted as observed district relocation. The primary package and alternative geographic products identify their respective boundaries, source licences and reporting masks. Cement works are placed at plant coordinates from the Global Cement and Concrete Tracker (CC BY 4.0) \citep{GEMCement2026}; coordinates derived from OpenStreetMap are used only in the separate ODbL companion.

Spatial closure is
\begin{equation}
 E_{p,t,c}^{\mathrm{national}}=\sum_g E_{p,t,c,g}^{\mathrm{placed}}
                      +E_{p,t,c}^{\mathrm{unallocated}},
 \label{eq:spatial}
\end{equation}
for species $p$, year $t$ and configuration $c$, on a compatible reporting boundary. Territories outside a statistical total are not silently included by a broader grid footprint. The maximum residual is reported at the product's actual numerical precision. Meaningful discrepancies are not absorbed into rounding.

The \SI{0.05}{\degree} grid stores native mass per cell and period. Flux density requires mass-unit conversion, the appropriate cell/country area and period duration. A country fraction already incorporated into allocated cell mass must not be applied again as though it were a flux. Atmospheric-model use additionally requires release height, speciation and source-boundary choices; these are not supplied merely by placing mass on a grid.

\section{Uncertainty and evaluation}
The annual conditional ensembles use common parameter identities, seeds and draw indices for greenhouse gases and air pollutants. Positive quantities use the specified mean-preserving distributions; physically bounded fractions use distributions matched to their stated conditional mean/variance where justified. The selected distributions are modelling assumptions. Their ensemble percentiles are not probabilities assigned to structural cases or empirical confidence statements about missing sources.

The nitrogen engine regenerates its material state and native N$_2$O/NH$_3$/NO$_x$ releases. Crop-residue burning sums crop-specific burned masses after applying their corresponding production draws, then applies the common combustion-completeness and species-factor draws; the same crop-activity identities reach the air calculation. The waste engine reruns its bounded and positive parameters and tests fate closure before accepting a realisation. Fugitives retain their source/activity and factor-group bounded model; identity-seeded streams preserve paired parameters across years. Coal heat and CO$_2$ factors are replayed by rank. Other streams retain the existing parameter graph and its documented approximation. Unsupported management/equipment shares and control performance are not independently perturbed into physically impossible partitions.

Fertiliser NH$_3$ factors retain their product-specific central values and share one table-wide multiplier in the conditional ensemble. This dependence assumption is not a measured cross-product covariance; the published interval is conditional on it.

Two annual dependence cases retain persistent parameters or resample annual activity errors. Trend intervals use paired realisations. Waste-management parameters that are held across the activity history persist in both cases. The four structural configurations remain separate; kiln operating seasons, fuel sulfur specifications and unresolved controls are additional deterministic cases without probabilities. Allocation uncertainty and unestimated sources are not folded into parameter percentiles.

The conditional level and paired-trend tables include native CO$_2$, CH$_4$ and N$_2$O masses and separate AR5 and AR6 aggregates. An analytical longitudinal sensitivity table, which does not replay every native material flow, is published with an explicit model-basis field and does not replace the conditional intervals. Pollutant intervals retain their chemical-mass convention. Driver products are calculated separately for greenhouse gases and each pollutant; correlation diagnostics identify sensitivity and are not summed as variance attribution. Prefix and reproducibility tests check that a shorter run reproduces the same initial realisations and that output-row ordering does not create new physical uncertainty.

Evaluation first matches period, gas/species, mass convention, category and fuel boundary. A half-weighted blend of adjacent calendar years is labelled a fiscal approximation unless monthly observations permit exact fiscal aggregation. Comparison with national reporting uses the stated IPCC/CRT crosswalk, not an unqualified join on category text. Comparisons with global inventories identify shared activity statistics and factors.

For pollutants, fossil-only versus fossil-plus-biomass industrial comparisons separate a coverage difference before testing factors. Power coal screening is not ranked alongside central admitted emissions. Provincial and district tests use common polygons and retain unallocated mass. Timing checks use source-month controls or evidence not used to create the same profile. Geographic comparisons distinguish construction proxies from independent evaluation evidence. Satellite columns or aerosol observations require appropriate transport, chemistry and retrieval treatment; a simple concentration--emission correlation is not used as validation.

\section{Data products, reproduction and access}
The FY2023--24 covered aggregate is \SI{\GHGReferenceTotal}{\mega\tonne} CO$_2$e on the AR5 basis. The same release gives \SI{\AirReferencePM}{\giga\gram} primary PM$_{2.5}$, \SI{\AirReferenceNOx}{\giga\gram} NO$_x$ and \SI{\AirReferenceNHthree}{\giga\gram} NH$_3$. The main article tabulates all core pollutants and shows selected native-gas maps, trends, structural cases and matched-boundary comparisons. These summaries are read from released products through a generated fact file.

The product catalogue identifies annual source records, administrative summaries, monthly activity/emissions, grids, coverage/factor/control evidence, conditional uncertainty, structural cases and evaluation tables. Source release versions and overlapping revisions remain traceable. Query interfaces use the same periods, units and canonical configurations as downloads and retain aliases for legacy identifiers where applicable.

Reproduction from compiled activity tables is distinguished from re-extraction of primary sources. The build records state which level was run, the dependencies, input hashes and outcome. The deposit includes the numerical products, source code, metadata, plotted data and reproduction instructions. Primary source documents are held by the project and redistributed only where their terms permit. Geographic alternatives derived from OpenStreetMap data under the Open Database Licence (ODbL) form a separate companion package.

\paragraph{Evidence tables behind the source checks.} Each parsed monthly source ships its own denominators and controls, and rows that fail a printed total are kept unrepaired in a quarantine file rather than corrected: \texttt{pbs\_\allowbreak minerals\_\allowbreak parse\_\allowbreak summary.csv} and \texttt{pbs\_\allowbreak minerals\_\allowbreak quarantine.csv} for PBS minerals, with \texttt{pakemis\_\allowbreak minerals\_\allowbreak panel\_\allowbreak screening.csv} recording which mineral-years of the province panel are whole; \texttt{pbs\_\allowbreak qim\_\allowbreak parse\_\allowbreak summary.csv} and \texttt{pbs\_\allowbreak qim\_\allowbreak quarantine.csv} for the Quantum Index of Manufacturing; \texttt{ocac\_\allowbreak parse\_\allowbreak summary.csv} and \texttt{ocac\_\allowbreak direct\_\allowbreak edition\_\allowbreak quarantine.csv} for petroleum sales, with \texttt{ocac\_\allowbreak edition\_\allowbreak currency.csv} naming the edition each month is built on and \texttt{ocac\_\allowbreak direct\_\allowbreak edition\_\allowbreak crosscheck.csv} comparing differenced months with OCAC's single-month editions; \texttt{nfdc\_\allowbreak monthly\_\allowbreak parse\_\allowbreak summary.csv} and \texttt{nfdc\_\allowbreak monthly\_\allowbreak quarantine.csv} for fertiliser offtake; and \texttt{apcma\_\allowbreak dispatch\_\allowbreak parse\_\allowbreak summary.csv} and \texttt{apcma\_\allowbreak dispatch\_\allowbreak quarantine.csv} for cement despatch. Monthly shapes are compared across sources in \texttt{limestone\_\allowbreak qim\_\allowbreak cement\_\allowbreak crosscheck.csv} and \texttt{apcma\_\allowbreak qim\_\allowbreak cement\_\allowbreak crosscheck.csv}, the residential temperature profile against ERA5 in \texttt{era5\_\allowbreak degree\_\allowbreak day\_\allowbreak crosscheck.csv} (timing, not level) and the rice calendar against FAO harvest windows in \texttt{fao\_\allowbreak crop\_\allowbreak calendar\_\allowbreak crosscheck.csv}. For coal, \texttt{pakemis\_\allowbreak coal\_\allowbreak allocation\_\allowbreak identity.csv} tests whether the consumer split is a measurement or a redistribution, \texttt{pakemis\_\allowbreak power\_\allowbreak coal\_\allowbreak energy\_\allowbreak adjudication.csv} asks whether NEPRA's aggregate coal energy and its per-station coal mass can describe the same fleet, \texttt{power\_\allowbreak table3\_\allowbreak hypotheses.csv} enumerates the candidate explanations without selecting one, and \texttt{power\_\allowbreak toe\_\allowbreak reconciliation.csv} lays out the NEPRA and K-Electric reconciliation by year and fuel. \texttt{pakemis\_\allowbreak gas\_\allowbreak flow\_\allowbreak crosswalk.csv} classifies every delivered gas stream by the kind of flow it is. Power-plant controls are recorded in \texttt{pakemis\_\allowbreak air\_\allowbreak power\_\allowbreak control\_\allowbreak evidence.csv} as public design evidence, without inferring demonstrated operation. \texttt{pakemis\_\allowbreak faostat\_\allowbreak nitrogen\_\allowbreak flows.csv} compares managed soils with FAOSTAT at the nitrogen flows, and \texttt{pakemis\_\allowbreak nitrogen\_\allowbreak partition\_\allowbreak sensitivity.csv} gives the greenhouse-gas consequences of the all-NH$_3$ chemical-attribution diagnostic, with no probability or bound. \texttt{pakemis\_\allowbreak wastewater\_\allowbreak method\_\allowbreak comparison.csv} sets domestic wastewater N$_2$O under the 2019 Refinement beside the 2006 single-equation method. \texttt{pakemis\_\allowbreak primary\_\allowbreak geography\_\allowbreak licence\_\allowbreak audit.csv} records each primary-geography dependency and the scope of the separate ODbL companion, and \texttt{pakemis\_\allowbreak citation\_\allowbreak deposit\_\allowbreak audit.csv} lists every output path the manuscripts cite with whether the deposit contains it.

A dataset DOI and external review-access deposit are pending. Neither this documentation nor the local package claims journal submission, external publication or acceptance.

\appendix
%% GENERATED by scripts/make_v3_product_list.py from data/catalog/artifact_roles.csv
%% (role=canonical), outputs/pakemis_product_catalog.csv and each file's leading comment.
%% Do not edit: run the script. --check compares bytes.
\section{Canonical data products}
\label{app:products}
The release ships 233 canonical products. Each is listed once below, by family, with the file name under \texttt{outputs/} and one line on what it holds. Descriptions are taken from the product catalogue, the file's own header or its declared role.

\subsection*{Semantic query layer}
\begin{itemize}\setlength{\itemsep}{1pt}\footnotesize
\item \texttt{semantic/\allowbreak dim\_\allowbreak activity\_\allowbreak columns.csv}: Semantic dimension table: activity columns.
\item \texttt{semantic/\allowbreak dim\_\allowbreak category.csv}: Semantic dimension table: category.
\item \texttt{semantic/\allowbreak dim\_\allowbreak configuration.csv}: Semantic dimension table: configuration.
\item \texttt{semantic/\allowbreak dim\_\allowbreak evidence\_\allowbreak class.csv}: Semantic dimension table: evidence class.
\item \texttt{semantic/\allowbreak dim\_\allowbreak gas.csv}: Semantic dimension table: gas.
\item \texttt{semantic/\allowbreak dim\_\allowbreak geography.csv}: Semantic dimension table: geography.
\item \texttt{semantic/\allowbreak dim\_\allowbreak gwp.csv}: Semantic dimension table: gwp.
\item \texttt{semantic/\allowbreak dim\_\allowbreak month.csv}: Semantic dimension table: month.
\item \texttt{semantic/\allowbreak dim\_\allowbreak notation\_\allowbreak key.csv}: Semantic dimension table: notation key.
\item \texttt{semantic/\allowbreak dim\_\allowbreak period.csv}: Dimension tables: the closed vocabularies every fact row keys into.
\item \texttt{semantic/\allowbreak dim\_\allowbreak pollutant.csv}: Semantic dimension table: pollutant.
\item \texttt{semantic/\allowbreak dim\_\allowbreak stream.csv}: Semantic dimension table: stream.
\item \texttt{semantic/\allowbreak fact\_\allowbreak activity\_\allowbreak aviation.csv}: Semantic fact table: activity aviation.
\item \texttt{semantic/\allowbreak fact\_\allowbreak activity\_\allowbreak crops.csv}: Semantic fact table: activity crops.
\item \texttt{semantic/\allowbreak fact\_\allowbreak activity\_\allowbreak fertiliser.csv}: Semantic fact table: activity fertiliser.
\item \texttt{semantic/\allowbreak fact\_\allowbreak activity\_\allowbreak fuel\_\allowbreak sales\_\allowbreak area.csv}: Semantic fact table: activity fuel sales area.
\item \texttt{semantic/\allowbreak fact\_\allowbreak activity\_\allowbreak fuel\_\allowbreak sales\_\allowbreak national.csv}: Semantic fact table: activity fuel sales national.
\item \texttt{semantic/\allowbreak fact\_\allowbreak activity\_\allowbreak gas\_\allowbreak network.csv}: Semantic fact table: activity gas network.
\item \texttt{semantic/\allowbreak fact\_\allowbreak activity\_\allowbreak gas\_\allowbreak supply.csv}: Semantic fact table: activity gas supply.
\item \texttt{semantic/\allowbreak fact\_\allowbreak activity\_\allowbreak household\_\allowbreak fuel.csv}: Semantic fact table: activity household fuel.
\item \texttt{semantic/\allowbreak fact\_\allowbreak activity\_\allowbreak industry\_\allowbreak fuel.csv}: Semantic fact table: activity industry fuel.
\item \texttt{semantic/\allowbreak fact\_\allowbreak activity\_\allowbreak industry\_\allowbreak production.csv}: Semantic fact table: activity industry production.
\item \texttt{semantic/\allowbreak fact\_\allowbreak activity\_\allowbreak livestock.csv}: Semantic fact table: activity livestock.
\item \texttt{semantic/\allowbreak fact\_\allowbreak activity\_\allowbreak power\_\allowbreak monthly.csv}: Semantic fact table: activity power monthly.
\item \texttt{semantic/\allowbreak fact\_\allowbreak activity\_\allowbreak power\_\allowbreak plants.csv}: Semantic fact table: activity power plants.
\item \texttt{semantic/\allowbreak fact\_\allowbreak activity\_\allowbreak railways.csv}: Semantic fact table: activity railways.
\item \texttt{semantic/\allowbreak fact\_\allowbreak activity\_\allowbreak source\_\allowbreak details.csv}: Semantic fact table: activity source details.
\item \texttt{semantic/\allowbreak fact\_\allowbreak activity\_\allowbreak vehicles.csv}: Semantic fact table: activity vehicles.
\item \texttt{semantic/\allowbreak fact\_\allowbreak activity\_\allowbreak waste\_\allowbreak flows.csv}: Semantic fact table: activity waste flows.
\item \texttt{semantic/\allowbreak fact\_\allowbreak activity\_\allowbreak waste\_\allowbreak services.csv}: Semantic fact table: activity waste services.
\item \texttt{semantic/\allowbreak fact\_\allowbreak adm1\_\allowbreak month.csv}: Semantic fact table: provincial (ADM1) month.
\item \texttt{semantic/\allowbreak fact\_\allowbreak adm1\_\allowbreak year.csv}: Semantic fact table: provincial (ADM1) year.
\item \texttt{semantic/\allowbreak fact\_\allowbreak adm2\_\allowbreak month.csv}: Semantic fact table: district (ADM2) month.
\item \texttt{semantic/\allowbreak fact\_\allowbreak adm2\_\allowbreak year.csv}: Semantic fact table: district (ADM2) year.
\item \texttt{semantic/\allowbreak fact\_\allowbreak air\_\allowbreak absence\_\allowbreak screening.csv}: Semantic fact table: air-pollutant absence screening.
\item \texttt{semantic/\allowbreak fact\_\allowbreak air\_\allowbreak adm1.csv}: Semantic fact table: air-pollutant provincial (ADM1).
\item \texttt{semantic/\allowbreak fact\_\allowbreak air\_\allowbreak adm1\_\allowbreak year.csv}: Semantic fact table: air-pollutant provincial (ADM1) year.
\item \texttt{semantic/\allowbreak fact\_\allowbreak air\_\allowbreak adm2.csv}: Semantic fact table: air-pollutant district (ADM2).
\item \texttt{semantic/\allowbreak fact\_\allowbreak air\_\allowbreak adm2\_\allowbreak year.csv}: Semantic fact table: air-pollutant district (ADM2) year.
\item \texttt{semantic/\allowbreak fact\_\allowbreak air\_\allowbreak annual.csv}: Semantic fact table: air-pollutant annual.
\item \texttt{semantic/\allowbreak fact\_\allowbreak air\_\allowbreak conditional.csv}: Semantic fact table: air-pollutant conditional.
\item \texttt{semantic/\allowbreak fact\_\allowbreak air\_\allowbreak lahore.csv}: Semantic fact table: air-pollutant lahore.
\item \texttt{semantic/\allowbreak fact\_\allowbreak air\_\allowbreak panel.csv}: Semantic fact table: air-pollutant panel.
\item \texttt{semantic/\allowbreak fact\_\allowbreak air\_\allowbreak source\_\allowbreak classes.csv}: Semantic fact table: air-pollutant source classes.
\item \texttt{semantic/\allowbreak fact\_\allowbreak air\_\allowbreak timeseries.csv}: Semantic fact table: air-pollutant timeseries.
\item \texttt{semantic/\allowbreak fact\_\allowbreak air\_\allowbreak trend.csv}: Semantic fact table: air-pollutant trend.
\item \texttt{semantic/\allowbreak fact\_\allowbreak annual.csv}: Typed fact tables (annual, calendar, monthly, ADM1, ADM2, unallocated, uncertainty, completeness, transfers, evidence) whose values are the products'.
\item \texttt{semantic/\allowbreak fact\_\allowbreak annual\_\allowbreak calendar.csv}: Semantic fact table: annual calendar.
\item \texttt{semantic/\allowbreak fact\_\allowbreak completeness.csv}: Semantic fact table: completeness.
\item \texttt{semantic/\allowbreak fact\_\allowbreak completeness\_\allowbreak mass.csv}: Semantic fact table: completeness mass.
\item \texttt{semantic/\allowbreak fact\_\allowbreak evidence.csv}: Semantic fact table: evidence.
\item \texttt{semantic/\allowbreak fact\_\allowbreak ghg\_\allowbreak conditional.csv}: Semantic fact table: greenhouse-gas conditional.
\item \texttt{semantic/\allowbreak fact\_\allowbreak ghg\_\allowbreak trend.csv}: Semantic fact table: greenhouse-gas trend.
\item \texttt{semantic/\allowbreak fact\_\allowbreak monthly.csv}: Semantic fact table: monthly.
\item \texttt{semantic/\allowbreak fact\_\allowbreak spatial\_\allowbreak basis\_\allowbreak year.csv}: Semantic fact table: spatial basis year.
\item \texttt{semantic/\allowbreak fact\_\allowbreak transfer.csv}: Semantic fact table: transfer.
\item \texttt{semantic/\allowbreak fact\_\allowbreak unallocated.csv}: Semantic fact table: unallocated.
\item \texttt{semantic/\allowbreak fact\_\allowbreak uncertainty.csv}: Semantic fact table: uncertainty.
\item \texttt{semantic/\allowbreak fact\_\allowbreak uncertainty\_\allowbreak series.csv}: Semantic fact table: uncertainty series.
\item \texttt{semantic/\allowbreak semantic\_\allowbreak release.json}: The semantic contract: release identity, the dimensions and their values, the defaults, the rounding and aggregation rules, the source digests.
\end{itemize}

\subsection*{Gridded products (NetCDF)}
\begin{itemize}\setlength{\itemsep}{1pt}\footnotesize
\item \texttt{pakemis\_\allowbreak air\_\allowbreak grid\_\allowbreak 0p05\_\allowbreak annual\_\allowbreak l\_\allowbreak ipcc\_\allowbreak fy2015\_\allowbreak 16\_\allowbreak fy2024\_\allowbreak 25.nc}: Annual and monthly covered air emissions.
\item \texttt{pakemis\_\allowbreak air\_\allowbreak grid\_\allowbreak 0p05\_\allowbreak annual\_\allowbreak l\_\allowbreak tan\_\allowbreak fy2015\_\allowbreak 16\_\allowbreak fy2024\_\allowbreak 25.nc}: Annual and monthly covered air emissions.
\item \texttt{pakemis\_\allowbreak air\_\allowbreak grid\_\allowbreak 0p05\_\allowbreak annual\_\allowbreak s\_\allowbreak ipcc\_\allowbreak fy2015\_\allowbreak 16\_\allowbreak fy2024\_\allowbreak 25.nc}: Annual and monthly covered air emissions.
\item \texttt{pakemis\_\allowbreak air\_\allowbreak grid\_\allowbreak 0p05\_\allowbreak annual\_\allowbreak s\_\allowbreak tan\_\allowbreak fy2015\_\allowbreak 16\_\allowbreak fy2024\_\allowbreak 25.nc}: Annual and monthly covered air emissions.
\item \texttt{pakemis\_\allowbreak air\_\allowbreak grid\_\allowbreak 0p05\_\allowbreak monthly\_\allowbreak l\_\allowbreak ipcc\_\allowbreak fy2015\_\allowbreak 16\_\allowbreak fy2024\_\allowbreak 25.nc}: Annual and monthly covered air emissions.
\item \texttt{pakemis\_\allowbreak air\_\allowbreak grid\_\allowbreak 0p05\_\allowbreak monthly\_\allowbreak l\_\allowbreak tan\_\allowbreak fy2015\_\allowbreak 16\_\allowbreak fy2024\_\allowbreak 25.nc}: Annual and monthly covered air emissions.
\item \texttt{pakemis\_\allowbreak air\_\allowbreak grid\_\allowbreak 0p05\_\allowbreak monthly\_\allowbreak s\_\allowbreak ipcc\_\allowbreak fy2015\_\allowbreak 16\_\allowbreak fy2024\_\allowbreak 25.nc}: Annual and monthly covered air emissions.
\item \texttt{pakemis\_\allowbreak air\_\allowbreak grid\_\allowbreak 0p05\_\allowbreak monthly\_\allowbreak s\_\allowbreak tan\_\allowbreak fy2015\_\allowbreak 16\_\allowbreak fy2024\_\allowbreak 25.nc}: Annual and monthly covered air emissions.
\item \texttt{pakemis\_\allowbreak ghg\_\allowbreak grid\_\allowbreak 0p05\_\allowbreak annual\_\allowbreak l\_\allowbreak ipcc\_\allowbreak fy2015\_\allowbreak 16\_\allowbreak fy2024\_\allowbreak 25.nc}: Annual native gas per cell for named structural configuration.
\item \texttt{pakemis\_\allowbreak ghg\_\allowbreak grid\_\allowbreak 0p05\_\allowbreak annual\_\allowbreak l\_\allowbreak tan\_\allowbreak fy2015\_\allowbreak 16\_\allowbreak fy2024\_\allowbreak 25.nc}: Annual native gas per cell for named structural configuration.
\item \texttt{pakemis\_\allowbreak ghg\_\allowbreak grid\_\allowbreak 0p05\_\allowbreak annual\_\allowbreak s\_\allowbreak ipcc\_\allowbreak fy2015\_\allowbreak 16\_\allowbreak fy2024\_\allowbreak 25.nc}: Annual native gas per cell for named structural configuration.
\item \texttt{pakemis\_\allowbreak grid\_\allowbreak 0p05\_\allowbreak annual\_\allowbreak ref\_\allowbreak fy2015\_\allowbreak 16\_\allowbreak fy2024\_\allowbreak 25.nc}: Ten fiscal periods at 0.05 degrees, three spatial views, three gases, in native mass.
\item \texttt{pakemis\_\allowbreak grid\_\allowbreak 0p05\_\allowbreak monthly.nc}: The 0.05 degree monthly cube in native gas mass, with district and provincial views, for the reference year.
\item \texttt{pakemis\_\allowbreak grid\_\allowbreak 0p05\_\allowbreak monthly\_\allowbreak ledger.nc}: Configuration L on the 0.05-degree monthly grid the three gas layers, CO2 bit-identical to the reference cube.
\item \texttt{pakemis\_\allowbreak grid\_\allowbreak 0p1\_\allowbreak edgar\_\allowbreak aligned.nc}: The annual 0.05-degree cube aggregated 2x2 at the odd parity that tiles EDGAR's 0.1-degree lattice exactly.
\end{itemize}

\subsection*{Uncertainty and structural cases}
\begin{itemize}\setlength{\itemsep}{1pt}\footnotesize
\item \texttt{pakemis\_\allowbreak air\_\allowbreak absence\_\allowbreak screening\_\allowbreak by\_\allowbreak year.csv}: What the panel's typed absences plausibly hold, per period and species.
\item \texttt{pakemis\_\allowbreak air\_\allowbreak conditional\_\allowbreak uncertainty.csv}: Conditional pollutant intervals with shared nitrogen, waste and combustion inputs.
\item \texttt{pakemis\_\allowbreak air\_\allowbreak structural\_\allowbreak sensitivities.csv}: Source-specific structural cases with no scenario probabilities.
\item \texttt{pakemis\_\allowbreak air\_\allowbreak trend\_\allowbreak uncertainty.csv}: Paired-realisations changes under declared temporal dependence.
\item \texttt{pakemis\_\allowbreak air\_\allowbreak uncertainty\_\allowbreak drivers.csv}: Conditional uncertainty sensitivity by pollutant.
\item \texttt{pakemis\_\allowbreak air\_\allowbreak uncertainty\_\allowbreak parameters.csv}: Parameters and constraints used in pollutant ensembles.
\item \texttt{pakemis\_\allowbreak air\_\allowbreak uncertainty\_\allowbreak split\_\allowbreak by\_\allowbreak year.csv}: The two uncertainty halves reported separately per (period, configuration, pollutant).
\item \texttt{pakemis\_\allowbreak bridge\_\allowbreak configurations.csv}: Records: bridge configurations.
\item \texttt{pakemis\_\allowbreak ensemble\_\allowbreak ledger\_\allowbreak summary.csv}: Configuration L on the same draw\_id as the reference ensemble.
\item \texttt{pakemis\_\allowbreak ensemble\_\allowbreak spatial\_\allowbreak conditionals.csv}: Spatial alternatives as conditional ensemble blocks.
\item \texttt{pakemis\_\allowbreak ensemble\_\allowbreak summary.csv}: A 10,000-draw prefix of the national Monte Carlo carried to months and provinces on one draw\_id, with both closure identities holding inside every.
\item \texttt{pakemis\_\allowbreak ghg\_\allowbreak conditional\_\allowbreak uncertainty.csv}: Configuration-conditional native gases and separate GWP aggregates.
\item \texttt{pakemis\_\allowbreak ghg\_\allowbreak trend\_\allowbreak uncertainty.csv}: Paired native-gas and separate GWP conditional uncertainty.
\item \texttt{pakemis\_\allowbreak ghg\_\allowbreak uncertainty\_\allowbreak drivers.csv}: Paired native-gas and separate GWP conditional uncertainty.
\item \texttt{pakemis\_\allowbreak ghg\_\allowbreak uncertainty\_\allowbreak parameters.csv}: Paired native-gas and separate GWP conditional uncertainty.
\item \texttt{pakemis\_\allowbreak joint\_\allowbreak uncertainty.csv}: Joint uncertainty by category and by the resolution of the figure being read.
\item \texttt{pakemis\_\allowbreak longitudinal\_\allowbreak prefix.csv}: Legacy analytical sensitivity.
\item \texttt{pakemis\_\allowbreak longitudinal\_\allowbreak uncertainty.csv}: Ten annual analytical bands under three temporal-dependence assumptions.
\item \texttt{pakemis\_\allowbreak structural\_\allowbreak trajectories.csv}: Every registered structural axis in every panel year.
\item \texttt{pakemis\_\allowbreak trend\_\allowbreak uncertainty.csv}: Draw-wise paired differences from the legacy analytical graph, restricted to sectors reported in both years.
\item \texttt{power\_\allowbreak air\_\allowbreak screening\_\allowbreak fy2015\_\allowbreak 16\_\allowbreak fy2024\_\allowbreak 25.csv}: Power's uncontrolled screening scenario for the same ten years, shipped beside the covered panel and never inside it.
\end{itemize}

\subsection*{Evaluation and comparison}
\begin{itemize}\setlength{\itemsep}{1pt}\footnotesize
\item \texttt{pakemis\_\allowbreak air\_\allowbreak lahore\_\allowbreak by\_\allowbreak year.csv}: Lahore district's air pollutants by sector and year.
\item \texttt{pakemis\_\allowbreak btr\_\allowbreak comparison.csv}: PakEmis against Pakistan's own national submission.
\item \texttt{pakemis\_\allowbreak builtup\_\allowbreak comparison\_\allowbreak adm2.csv}: Records: builtup comparison district (ADM2).
\item \texttt{pakemis\_\allowbreak comparison\_\allowbreak crosswalk.csv}: What each external comparison compares.
\item \texttt{pakemis\_\allowbreak enteric\_\allowbreak btr\_\allowbreak bridge.csv}: Enteric CH4 by species, this inventory (FY2020-21, REF\_2023\_24, AR5) against Pakistan's.
\item \texttt{pakemis\_\allowbreak industry\_\allowbreak edgar\_\allowbreak divergence.csv}: The engine's industrial series against EDGAR'S where A fiscal blend exists.
\item \texttt{pakemis\_\allowbreak validation\_\allowbreak catalog.csv}: Per-catalog-row validation status.
\item \texttt{pakemis\_\allowbreak validation\_\allowbreak five\_\allowbreak classes.csv}: Every validation-catalogue entry filed under what the comparison can establish.
\item \texttt{pakemis\_\allowbreak validation\_\allowbreak five\_\allowbreak classes\_\allowbreak summary.csv}: Counts of pakemis\_validation\_five\_classes.csv by class.
\item \texttt{pakemis\_\allowbreak validation\_\allowbreak measured.csv}: Measured validation, one row per check records each check read, how many it.
\item \texttt{pakemis\_\allowbreak validation\_\allowbreak scorecard.csv}: The validation scorecard the (property x grain) space, one row.
\item \texttt{unsd\_\allowbreak coal\_\allowbreak supply\_\allowbreak control.csv}: UNSD energy balances - Pakistan primary coal and peat supply block, tj, per edition x calendar year.
\item \texttt{unsd\_\allowbreak industry\_\allowbreak combustion\_\allowbreak control.csv}: Independent 1.A.2 industrial-combustion control total from the UNSD.
\end{itemize}

\subsection*{Administrative and spatial records}
\begin{itemize}\setlength{\itemsep}{1pt}\footnotesize
\item \texttt{air\_\allowbreak pollutants\_\allowbreak adm1.csv}: Legacy FY2023-24 provincial air layer of seven sectors.
\item \texttt{air\_\allowbreak pollutants\_\allowbreak adm2.csv}: Legacy FY2023-24 district air layer of seven sectors (power rows are screening, labelled by estimate\_kind).
\item \texttt{pakemis\_\allowbreak adm1\_\allowbreak by\_\allowbreak year.csv}: Ten fiscal years at ADM1 and ADM2, annually and monthly, under four configurations.
\item \texttt{pakemis\_\allowbreak adm1\_\allowbreak ledger\_\allowbreak 2024.csv}: Configuration L at ADM1 grain, stream x gas x month x province.
\item \texttt{pakemis\_\allowbreak adm2\_\allowbreak by\_\allowbreak year.csv}: Records: district (ADM2) by year.
\item \texttt{pakemis\_\allowbreak adm2\_\allowbreak crosswalk.csv}: Records: district (ADM2) crosswalk.
\item \texttt{pakemis\_\allowbreak adm2\_\allowbreak frame\_\allowbreak 2022.csv}: The OCHA COD-AB 2022 ADM2 frame (160 units, CC BY-IGO) crosswalked to the canonical 152-unit set.
\item \texttt{pakemis\_\allowbreak adm2\_\allowbreak ledger\_\allowbreak 2024.csv}: Configuration L at ADM2 grain, stream x gas x month x shape v2).
\item \texttt{pakemis\_\allowbreak adm2\_\allowbreak monthly.csv}: District monthly greenhouse-gas records for the reference year, without intervals.
\item \texttt{pakemis\_\allowbreak adm2\_\allowbreak monthly\_\allowbreak by\_\allowbreak year.csv}: Records: district (ADM2) monthly by year.
\item \texttt{pakemis\_\allowbreak adm2\_\allowbreak monthly\_\allowbreak sector.csv}: District monthly greenhouse-gas records by sector for the reference year, without intervals.
\item \texttt{pakemis\_\allowbreak adm2\_\allowbreak nesting\_\allowbreak check.csv}: Check that district agricultural allocations sum to their provincial controls.
\item \texttt{pakemis\_\allowbreak adm2\_\allowbreak sector\_\allowbreak by\_\allowbreak year.csv}: Records: district (ADM2) sector by year.
\item \texttt{pakemis\_\allowbreak adm2\_\allowbreak stream\_\allowbreak by\_\allowbreak year.csv}: The per-stream district vector for all ten periods.
\item \texttt{pakemis\_\allowbreak adm2\_\allowbreak stream\_\allowbreak reference.csv}: The reference configuration's per-stream ADM2 vectors, annual (Gg), all modules.
\item \texttt{pakemis\_\allowbreak air\_\allowbreak adm1\_\allowbreak by\_\allowbreak year.csv}: Ten-year provincial air totals at anchor-stream grain, the cube's own province marginals.
\item \texttt{pakemis\_\allowbreak air\_\allowbreak adm2\_\allowbreak by\_\allowbreak year.csv}: Ten-year district air totals on the canonical shape frame, aggregated per (configuration, sector, pollutant).
\item \texttt{pakemis\_\allowbreak air\_\allowbreak monthly\_\allowbreak unallocated.csv}: Records: air-pollutant monthly unallocated.
\item \texttt{pakemis\_\allowbreak air\_\allowbreak spatial\_\allowbreak support\_\allowbreak by\_\allowbreak year.csv}: Every air (period, configuration, anchor, pollutant) mass beside its weight-registry row.
\item \texttt{pakemis\_\allowbreak air\_\allowbreak unallocated\_\allowbreak by\_\allowbreak year.csv}: The air mass no geometry could place below ADM1, typed per period and configuration (the aviation cruise class).
\item \texttt{pakemis\_\allowbreak ghg\_\allowbreak configuration\_\allowbreak adm1.csv}: Four configurations placed with own component masses.
\item \texttt{pakemis\_\allowbreak ghg\_\allowbreak configuration\_\allowbreak adm1\_\allowbreak monthly.csv}: Four configurations placed with own component masses.
\item \texttt{pakemis\_\allowbreak ghg\_\allowbreak configuration\_\allowbreak adm2.csv}: All four configurations at ADM2.
\item \texttt{pakemis\_\allowbreak ghg\_\allowbreak configuration\_\allowbreak adm2\_\allowbreak monthly.csv}: Four configurations placed with own component masses.
\item \texttt{pakemis\_\allowbreak ghg\_\allowbreak configuration\_\allowbreak unallocated.csv}: Unallocated native and GWP mass for each configuration and fiscal month.
\item \texttt{pakemis\_\allowbreak grid\_\allowbreak fallbacks.csv}: Records: grid fallbacks.
\item \texttt{pakemis\_\allowbreak grid\_\allowbreak fallbacks\_\allowbreak by\_\allowbreak year.csv}: Records: grid fallbacks by year.
\item \texttt{pakemis\_\allowbreak grid\_\allowbreak provincial\_\allowbreak check.csv}: Records: grid provincial check.
\item \texttt{pakemis\_\allowbreak grid\_\allowbreak stream\_\allowbreak surfaces\_\allowbreak by\_\allowbreak year.csv}: Records: grid stream surfaces by year.
\item \texttt{pakemis\_\allowbreak kiln\_\allowbreak allocation\_\allowbreak adm2.csv}: Records: kiln allocation district (ADM2).
\item \texttt{pakemis\_\allowbreak power\_\allowbreak grid\_\allowbreak basis\_\allowbreak change.csv}: Comparison of the power geography under the alternative allocation bases.
\item \texttt{pakemis\_\allowbreak road\_\allowbreak length\_\allowbreak adm2.csv}: Records: road length district (ADM2).
\item \texttt{pakemis\_\allowbreak road\_\allowbreak links.csv}: Records: road links.
\item \texttt{pakemis\_\allowbreak road\_\allowbreak technology\_\allowbreak case.csv}: Road exhaust PM2.5 and bc under technology-specific tier 2 factors the central.
\item \texttt{pakemis\_\allowbreak road\_\allowbreak weights\_\allowbreak adm2.csv}: Records: road weights district (ADM2).
\item \texttt{pakemis\_\allowbreak spatial\_\allowbreak basis\_\allowbreak classes.csv}: The spatial evidence vocabulary (spatial\_time\_basis), defined in plain words.
\item \texttt{pakemis\_\allowbreak spatial\_\allowbreak basis\_\allowbreak shares.csv}: Mass-weighted share of each period's emissions by spatial evidence class, national.
\item \texttt{pakemis\_\allowbreak spatial\_\allowbreak basis\_\allowbreak shares\_\allowbreak by\_\allowbreak sector.csv}: Mass-weighted share of each period's emissions by spatial evidence class, per IPCC reporting sector.
\item \texttt{pakemis\_\allowbreak spatial\_\allowbreak brackets.csv}: Records: spatial brackets.
\item \texttt{pakemis\_\allowbreak spatial\_\allowbreak brackets\_\allowbreak adm2.csv}: Records: spatial brackets district (ADM2).
\item \texttt{pakemis\_\allowbreak spatial\_\allowbreak decomposition.csv}: Records: spatial decomposition.
\item \texttt{pakemis\_\allowbreak spatial\_\allowbreak exceptions.csv}: Allocations where the grid could not follow the declared spatial basis.
\item \texttt{pakemis\_\allowbreak spatial\_\allowbreak unallocated.csv}: The mass with no defensible placement below ADM1, retained rather than spread.
\item \texttt{pakemis\_\allowbreak spatial\_\allowbreak unallocated\_\allowbreak by\_\allowbreak year.csv}: Records: spatial unallocated by year.
\item \texttt{pakemis\_\allowbreak spatial\_\allowbreak weights\_\allowbreak by\_\allowbreak year.csv}: Records: spatial weights by year.
\item \texttt{pakemis\_\allowbreak stream\_\allowbreak support\_\allowbreak ledger.csv}: One row per (period, stream).
\item \texttt{pakemis\_\allowbreak unallocated\_\allowbreak activity\_\allowbreak controls.csv}: Records: unallocated activity controls.
\end{itemize}

\subsection*{Monthly records}
\begin{itemize}\setlength{\itemsep}{1pt}\footnotesize
\item \texttt{pakemis\_\allowbreak agriculture\_\allowbreak crop\_\allowbreak calendar.csv}: Which FAOSTAT calendar year each fiscal crop cell actually is.
\item \texttt{pakemis\_\allowbreak air\_\allowbreak monthly.csv}: Process-specific pollutant profiles with native mass and temporal evidence.
\item \texttt{pakemis\_\allowbreak air\_\allowbreak temporal\_\allowbreak support.csv}: Pollutant-specific temporal evidence shares.
\item \texttt{pakemis\_\allowbreak fy2025\_\allowbreak 26\_\allowbreak provisional\_\allowbreak monthly.csv}: Separate provisional source activity and emissions.
\item \texttt{pakemis\_\allowbreak ghg\_\allowbreak configuration\_\allowbreak monthly.csv}: Twelve process allocation months per source and fiscal year.
\item \texttt{pakemis\_\allowbreak monthly\_\allowbreak controls.csv}: Records: monthly controls.
\item \texttt{pakemis\_\allowbreak monthly\_\allowbreak coverage.csv}: What each month's timing rests on, one row per (month, sector, category, class).
\item \texttt{pakemis\_\allowbreak monthly\_\allowbreak coverage\_\allowbreak summary.csv}: Exactly 120 rows: what each month's timing rests on, and which sectors a month does not have at all.
\item \texttt{pakemis\_\allowbreak monthly\_\allowbreak data\_\allowbreak requests.csv}: Records: monthly data requests.
\item \texttt{pakemis\_\allowbreak monthly\_\allowbreak evidence.csv}: Records: monthly evidence.
\item \texttt{pakemis\_\allowbreak monthly\_\allowbreak fy2015\_\allowbreak 16\_\allowbreak fy2024\_\allowbreak 25.csv}: 120 months, July 2015 to June 2025, every non-absent component resolved into twelve months closing exactly to its annual figure.
\item \texttt{pakemis\_\allowbreak temporal\_\allowbreak decomposition.csv}: Records: temporal decomposition.
\item \texttt{pakemis\_\allowbreak temporal\_\allowbreak profiles.csv}: Records: temporal profiles.
\end{itemize}

\subsection*{Coverage and evidence registers}
\begin{itemize}\setlength{\itemsep}{1pt}\footnotesize
\item \texttt{completeness.csv}: The notation-key completeness table, one row per sector-category.
\item \texttt{pakemis\_\allowbreak air\_\allowbreak coverage\_\allowbreak matrix.csv}: The 700-cell coverage matrix (10 periods x 10 modules x 7 core species) under the closed IPCC notation keys.
\item \texttt{pakemis\_\allowbreak air\_\allowbreak source\_\allowbreak coverage.csv}: Process by species by period coverage.
\item \texttt{pakemis\_\allowbreak artifact\_\allowbreak inventory.csv}: Every file this release ships, with the role it declares and who names it.
\item \texttt{pakemis\_\allowbreak completeness\_\allowbreak mass.csv}: Unestimated streams joined to available activity-based screening quantities, with overlap and scenario assumptions retained.
\item \texttt{pakemis\_\allowbreak coverage\_\allowbreak matrix.csv}: Records: coverage matrix.
\item \texttt{pakemis\_\allowbreak coverage\_\allowbreak summary.csv}: Records: coverage summary.
\item \texttt{pakemis\_\allowbreak national\_\allowbreak series\_\allowbreak coverage.csv}: Records: national series coverage.
\item \texttt{pakemis\_\allowbreak product\_\allowbreak catalog.csv}: One row per logical product, not per file.
\item \texttt{pakemis\_\allowbreak product\_\allowbreak coverage.csv}: Records: product coverage.
\item \texttt{pakemis\_\allowbreak product\_\allowbreak facts.csv}: One authoritative answer per fact a manuscript states about the shape of this release.
\end{itemize}

\subsection*{Activity accounts, source controls and extension}
\begin{itemize}\setlength{\itemsep}{1pt}\footnotesize
\item \texttt{ogra\_\allowbreak gas\_\allowbreak class\_\allowbreak series.csv}: OGRA's sector-wise gas classes per fiscal year (MMCFD, average daily).
\item \texttt{pakemis\_\allowbreak agriculture\_\allowbreak basis\_\allowbreak change.csv}: Source/basis change from archived FAOSTAT-level/calendar activity to current fiscal crops/nutrients and adopted-anchor livestock reconstruction.
\item \texttt{pakemis\_\allowbreak agriculture\_\allowbreak fiscal\_\allowbreak series.csv}: Agriculture fiscal activity, FY2015-16 to FY2024-25.
\item \texttt{pakemis\_\allowbreak ammonia\_\allowbreak nitric\_\allowbreak series.csv}: Ammonia (2.B.1) and nitric acid (2.B.2) as nine-year series.
\item \texttt{pakemis\_\allowbreak brick\_\allowbreak activity\_\allowbreak reconstruction.csv}: Brick-kiln coal from A kiln count and A measured fuel rate.
\item \texttt{pakemis\_\allowbreak brick\_\allowbreak coal\_\allowbreak alternatives.csv}: Brick-kiln coal - the bounded alternatives, per fiscal year.
\item \texttt{pakemis\_\allowbreak cement\_\allowbreak activity\_\allowbreak sources.csv}: Records: cement activity sources.
\item \texttt{pakemis\_\allowbreak cement\_\allowbreak kiln\_\allowbreak coal\_\allowbreak reconstruction.csv}: Cement kiln coal reconstructed from clinker Q = K x H x s / NCV as a band per fiscal year.
\item \texttt{pakemis\_\allowbreak coal\_\allowbreak allocation\_\allowbreak cases.csv}: Dated conditional coal allocation cases.
\item \texttt{pakemis\_\allowbreak coal\_\allowbreak balance.csv}: National coal balance per fiscal year supply lanes vs HDIP's allocation vs NEPRA's.
\item \texttt{pakemis\_\allowbreak coal\_\allowbreak change\_\allowbreak decomposition.csv}: Coal emissions, year on year.
\item \texttt{pakemis\_\allowbreak coal\_\allowbreak source\_\allowbreak reconciliation.csv}: Industrial coal - source reconciliation one row per fiscal year x HDIP consumer.
\item \texttt{pakemis\_\allowbreak coal\_\allowbreak vintage\_\allowbreak ledger.csv}: Consecutive prints of the same reference period, paired.
\item \texttt{pakemis\_\allowbreak extension\_\allowbreak activity\_\allowbreak revisions.csv}: Source availability, revisions and controls for common core and candidate extension.
\item \texttt{pakemis\_\allowbreak extension\_\allowbreak source\_\allowbreak availability.csv}: Source availability, revisions and controls for common core and candidate extension.
\item \texttt{pakemis\_\allowbreak extension\_\allowbreak source\_\allowbreak controls.csv}: Source availability, revisions and controls for common core and candidate extension.
\item \texttt{pakemis\_\allowbreak fuel\_\allowbreak sulfur\_\allowbreak decision.csv}: The fuel-sulfur decision record (item 3), under the same acceptance policy as.
\item \texttt{pakemis\_\allowbreak fy2025\_\allowbreak 26\_\allowbreak provisional\_\allowbreak activity.csv}: Separate provisional source activity and emissions.
\item \texttt{pakemis\_\allowbreak fy2025\_\allowbreak 26\_\allowbreak provisional\_\allowbreak emissions.csv}: FY2025-26 source estimates with explicit partial coverage.
\item \texttt{pakemis\_\allowbreak hfc\_\allowbreak bank\_\allowbreak series.csv}: HFC emissions from the household equipment bank, FY2015-16..FY2023-24.
\item \texttt{pakemis\_\allowbreak industrial\_\allowbreak change\_\allowbreak attribution.csv}: Where the 1.A.2 decline comes from.
\item \texttt{pakemis\_\allowbreak industrial\_\allowbreak coal\_\allowbreak bridge.csv}: Industrial coal - the replacement bridge and reference decision per fiscal year x component.
\item \texttt{pakemis\_\allowbreak industrial\_\allowbreak coal\_\allowbreak decision.csv}: The decision record under the acceptance policy.
\item \texttt{pakemis\_\allowbreak industrial\_\allowbreak screening\_\allowbreak series.csv}: The remaining industrial categories as nine-year screening series.
\item \texttt{pakemis\_\allowbreak industry\_\allowbreak boundary\_\allowbreak ledger.csv}: The industrial boundary ledger.
\item \texttt{pakemis\_\allowbreak minerals\_\allowbreak controls.csv}: Independent source controls checked against final exported monthly/provincial panel.
\item \texttt{pakemis\_\allowbreak ocac\_\allowbreak extension\_\allowbreak revisions.csv}: Later publication overlap and unresolved sector by product revisions.
\item \texttt{pakemis\_\allowbreak power\_\allowbreak extension\_\allowbreak controls.csv}: Physical station versus system fuel and generation controls.
\item \texttt{pakemis\_\allowbreak rice\_\allowbreak geography\_\allowbreak crosscheck.csv}: The declared MapSPAM rice surface measured against the own-year GCD-Rice 2023 30 m classification, as within-province ADM2 shares with per-unit.
\item \texttt{pakemis\_\allowbreak series\_\allowbreak splice\_\allowbreak agriculture.csv}: Records: series splice agriculture.
\item \texttt{pakemis\_\allowbreak sulfur\_\allowbreak adoption\_\allowbreak effect.csv}: What the sulfur adoption moved, per stream and per year, with every sign shown.
\end{itemize}

\subsection*{Annual records}
\begin{itemize}\setlength{\itemsep}{1pt}\footnotesize
\item \texttt{national\_\allowbreak all\_\allowbreak sectors.csv}: Ten sectors on one reporting period, exact mass closure and executable cross-boundary ledgers.
\item \texttt{national\_\allowbreak all\_\allowbreak sectors\_\allowbreak ledger.csv}: Records: national all sectors ledger.
\item \texttt{pakemis\_\allowbreak air\_\allowbreak annual\_\allowbreak fy2015\_\allowbreak 16\_\allowbreak fy2024\_\allowbreak 25.csv}: Seven core pollutants in native mass with narrower additional species.
\item \texttt{pakemis\_\allowbreak air\_\allowbreak atomic\_\allowbreak annual.csv}: The ten-fiscal-year air panel's atoms.
\item \texttt{pakemis\_\allowbreak air\_\allowbreak power\_\allowbreak omission.csv}: What the covered national total leaves out at power stacks, per period and species.
\item \texttt{pakemis\_\allowbreak air\_\allowbreak source\_\allowbreak classes\_\allowbreak fy2015\_\allowbreak 16\_\allowbreak fy2024\_\allowbreak 25.csv}: The source-class air layer (industry, ippu, fugitives) at class grain, ten years.
\item \texttt{pakemis\_\allowbreak annual\_\allowbreak fy2015\_\allowbreak 16\_\allowbreak fy2024\_\allowbreak 25.csv}: Sector, category and gas for all 100 sector-years.
\item \texttt{pakemis\_\allowbreak atomic\_\allowbreak annual.csv}: The component grain beneath the panel.
\item \texttt{pakemis\_\allowbreak atomic\_\allowbreak annual\_\allowbreak aggregates.csv}: Records: atomic annual aggregates.
\item \texttt{pakemis\_\allowbreak claim\_\allowbreak matrix.csv}: Records: claim matrix.
\item \texttt{pakemis\_\allowbreak deposit\_\allowbreak manifest.csv}: Every file the archive deposit ships, with the licence under which it ships.
\item \texttt{pakemis\_\allowbreak dung\_\allowbreak residual\_\allowbreak ledger.csv}: The unresolved dung residual (configuration S) by species.
\item \texttt{pakemis\_\allowbreak facility\_\allowbreak registry.csv}: Records: facility registry.
\item \texttt{pakemis\_\allowbreak flatness\_\allowbreak crosswalk.csv}: The crosswalk between the TE-class flat share and the physical flatness decomposition.
\item \texttt{pakemis\_\allowbreak ghg\_\allowbreak configuration\_\allowbreak atoms.csv}: Native GHG emissions for S\_TAN, L\_TAN, S\_IPCC and L\_IPCC.
\item \texttt{pakemis\_\allowbreak national\_\allowbreak series.csv}: Records: national series.
\item \texttt{pakemis\_\allowbreak nitrogen\_\allowbreak no\_\allowbreak pathways.csv}: Agricultural NOx by pathway.
\item \texttt{pakemis\_\allowbreak paper\_\allowbreak artefact\_\allowbreak hashes.csv}: Records: paper artefact hashes.
\item \texttt{pakemis\_\allowbreak power\_\allowbreak profile\_\allowbreak variant.csv}: Records: power profile variant.
\item \texttt{pakemis\_\allowbreak power\_\allowbreak station\_\allowbreak status.csv}: Coal stations carrying each period's coal geography, joined to GEM unit start.
\item \texttt{pakemis\_\allowbreak source\_\allowbreak register.csv}: Every registered source: whether the deposit ships it, and if not, how to obtain it.
\item \texttt{pakemis\_\allowbreak source\_\allowbreak streams.csv}: Records: source streams.
\end{itemize}



%% GENERATED by scripts/make_release_cover.py from outputs/release_manifest.json.
%% Do not edit: the edit will be overwritten and CI diffs the result.
\newcommand{\ReleaseCover}{%
\begin{center}\small
\setlength{\fboxsep}{4pt}%
\fbox{\begin{minipage}{\dimexpr\linewidth-2\fboxsep-2\fboxrule\relax}\footnotesize
\textbf{Release provenance.} This document describes PakEmis under the period, configuration and source boundaries specified in its methods. The identifiers below link it to the accompanying products and source register.\\[2pt]
\begin{tabular}{@{}p{0.32\linewidth}@{\hspace{0.03\linewidth}}p{0.65\linewidth}@{}}
Release version & \texttt{1.0.0rc69} \\
Commit & \texttt{d8d874622f9a} \\
Reporting period & July 2015--June 2025 (ten fiscal years) (2015-07-01 to 2025-06-30) \\
Reference & S\_TAN (household survey; TAN nitrogen) \\
FY2023--24 total & 444.61 Mt \COe{} (AR5) under S\_TAN, the reference; 449.85 Mt under L\_TAN \\
Release manifest & \texttt{ac0ffe93ad9fcd88} (SHA-256 of \texttt{outputs/release\_manifest.json}, first 16) \\
Source manifest & \texttt{c1be6d314cfce1ea} (SHA-256 of \texttt{data/manifests/source\_manifest.csv}, first 16) \\
Data-description article & \texttt{paper\_national/} (ESSD, in preparation), same release
\end{tabular}
\end{minipage}}
\end{center}}

\ReleaseCover

\begin{thebibliography}{99}

\bibitem[Crippa et~al.(2020)]{Crippa2020}
Crippa, M., Solazzo, E., Huang, G., Guizzardi, D., Koffi, E., Muntean, M., Schieberle, C., Friedrich, R., and Janssens-Maenhout, G.: High resolution temporal profiles in the Emissions Database for Global Atmospheric Research, Scientific Data, 7, 121, \doi{10.1038/s41597-020-0462-2}, 2020.

\bibitem[FAO(2025)]{FAOSTAT2025}
FAO: FAOSTAT: emissions totals, crops and livestock, fertilisers by nutrient, and forestry production and trade, Food and Agriculture Organization of the United Nations, Rome [data sets], \url{https://www.fao.org/faostat/}, 2025 snapshot used in the source inventory; catalogue checked 30 September 2026, 2025.

\bibitem[Global Energy Monitor(2026)]{GEMCement2026}
Global Energy Monitor: Global Cement and Concrete Tracker, plant-level data, July 2026 (standard copy V1), Global Energy Monitor [data set], licensed under CC BY 4.0, 2026.

\bibitem[IPCC(2006)]{IPCC2006}
IPCC: 2006 IPCC Guidelines for National Greenhouse Gas Inventories, edited by: Eggleston, H. S., Buendia, L., Miwa, K., Ngara, T., and Tanabe, K., Institute for Global Environmental Strategies, Hayama, Japan, \url{https://www.ipcc-nggip.iges.or.jp/public/2006gl/}, 2006.

\bibitem[IPCC(2013)]{IPCCAR5}
IPCC: Climate Change 2013: The Physical Science Basis, Cambridge University Press, Cambridge, \doi{10.1017/CBO9781107415324}, 2013.

\bibitem[IPCC(2019)]{IPCC2019}
IPCC: 2019 Refinement to the 2006 IPCC Guidelines for National Greenhouse Gas Inventories, edited by: Calvo Buendia, E., Tanabe, K., Kranjc, A., Baasansuren, J., Fukuda, M., Ngarize, S., Osako, A., Pyrozhenko, Y., Shermanau, P., and Federici, S., IPCC, Switzerland, \url{https://www.ipcc-nggip.iges.or.jp/public/2019rf/}, 2019.

\bibitem[IPCC(2021)]{IPCCAR6}
IPCC: Climate Change 2021: The Physical Science Basis, Cambridge University Press, Cambridge, \doi{10.1017/9781009157896}, 2021.

\bibitem[Kaza et~al.(2018)]{Kaza2018}
Kaza, S., Yao, L., Bhada-Tata, P., and Van Woerden, F.: What a Waste 2.0: A Global Snapshot of Solid Waste Management to 2050, World Bank, Washington, DC, \doi{10.1596/978-1-4648-1329-0}, 2018.

\bibitem[Ministry of Finance(2024)]{MoF2024}
Ministry of Finance: Pakistan Economic Survey 2023--24, Government of Pakistan, Islamabad, chapter 2, Agriculture, 2024.

\bibitem[Ministry of Finance(2026)]{MoF2026}
Ministry of Finance: Pakistan Economic Survey 2025--26, Government of Pakistan, Islamabad, agriculture chapter and statistical annexes, \url{https://www.finance.gov.pk/survey/chapter_26/2_Agriculture.pdf}, 2026.

\bibitem[NEPRA(2024)]{NEPRA2024}
NEPRA: State of Industry Report 2024, National Electric Power Regulatory Authority, Islamabad, \url{https://nepra.org.pk/publications/State%20of%20Industry%20Reports/State%20of%20Industry%20Report%202024.pdf}, 2024.

\bibitem[NFDC(2024)]{NFDC2024}
NFDC: Fertilizer Review 2023--24, Publication No. 2/2024, National Fertilizer Development Centre, Government of Pakistan, Islamabad, \url{https://nfdc.gov.pk/publications/Annual%20Fertilizer%20Review%202023-24.pdf}, 2024.

\bibitem[NFDC(2025)]{NFDC2025}
NFDC: Annual Fertilizer Review 2024--25, National Fertilizer Development Centre, Government of Pakistan, \url{https://mnfsr.gov.pk/SiteImage/Misc/files/AFR2024-25.pdf}, 2025.

\bibitem[Pak-EPA and EPMC(1996)]{PakEPA1996}
Pak-EPA and EPMC: National Study on Privatization of Solid Waste Management in Eight Cities of Pakistan, Pakistan Environmental Protection Agency, Islamabad, 1996.

\bibitem[PBS(2019)]{HIES2019}
PBS: Household Integrated Economic Survey 2018--19, Pakistan Bureau of Statistics, Islamabad, survey report and microdata, 2019.

\bibitem[PBS(2023)]{PBSCensus2023}
PBS: Seventh Population and Housing Census 2023, Pakistan Bureau of Statistics, Islamabad, census results, 2023.

\bibitem[PBS(2026a)]{PBSAgri2026}
PBS: Agricultural Census 2024, Pakistan Bureau of Statistics, Islamabad, report collection, \url{https://www.pbs.gov.pk/agriculture-census/}, accessed 30 September 2026, 2026a.

\bibitem[PBS(2026b)]{PBSHIES2026}
PBS: Household Integrated Economic Survey (HIES) 2024--25, Pakistan Bureau of Statistics, Islamabad, \url{https://www.pbs.gov.pk/household-integrated-economic-survey-hies-2024-25/}, accessed 30 September 2026, 2026b.

\bibitem[WHO and UNICEF(2023)]{JMP2023}
WHO and UNICEF: Joint Monitoring Programme for Water Supply, Sanitation and Hygiene: Pakistan household country file, 2023 snapshot, \url{https://washdata.org/data/downloads}, 2023.


\end{thebibliography}
\end{document}
